% Auto-generated by emit_k_of_d_corruption_table.py (do not edit by hand)
\begin{tabular}{p{0.16\textwidth}p{0.35\textwidth}p{0.38\textwidth}}
\toprule
\textbf{Claim} & \textbf{Current evidence} & \textbf{Boundary still open} \\
\midrule
T1 quality-blind impossibility & Architectural motivation: score-only fusion cannot disambiguate identical score vectors with different latent reliability states. & Formal example is theoretical; no separate empirical table is needed beyond collapse construction. \\
T2 global-KS mixture confounding & MVTec category-aware reference experiments show global references can over-activate under heterogeneous categories. & Needs a pure controlled mixture-shift table with no within-category corruption. \\
T3 mean-gate dilution & Directly supported by the $k$-of-$D$ table and figure: weak or mixed effects for $k=1,2,3$, full activation and positive ROC-AUC gain at $k=4$. & The ideal theorem assumes failed-domain reliability is zero; the implemented estimator keeps minimum reliability above $0.34$. \\
T4 risk dominance & The theorem gives measurable terms $(q_0,q_1,\Delta_0,\Delta_1)$ for when switching should win. & These terms are not yet reported as a deployment prevalence sensitivity table. \\
T5 finite-sample certificate & Certificate utility exists and defines fired-case paired benefit. & Needs an audited deployment-window LCB table before making operational safety claims. \\
T6 KS false-fire/power & KS is the empirically necessary trigger in component ablation. & Needs sample-size and multiple-testing power curves before claiming calibrated drift detection. \\
\bottomrule
\end{tabular}
